\begin{proposition}[an exact spin-flip equation forces the conference shell]
\label{prop:kramers-conference-shell}
Let $J=J_7$. Choose a distinguished qubit $s$
and local $Z$ axes, and define
\[
 H_s=-\frac12Z_s-\frac14\sum_{i\ne s}Z_i,
 \qquad {\cal A}_s=H_sJ.
\]
A normalised pure seven-qubit state satisfying
$z={\cal A}_sz$ lies in the \emph{complex
pair/complement family} of
Theorem~\ref{thm:complexsos},
up to a special-qubit diagonal $SU(2)$ phase.
Conversely all members of that family
satisfy this equation in that phase frame.
Thus every such $z$ has $\pi_3(z)\ge121/917$, with equality
only on the icosahedral LU/permutation orbit.
In its displayed conference frame the
icosahedral state satisfies
$\Psiico=H_sJ\Psiico$.
\end{proposition}
\begin{proof}
Since $J^2=-I$ and $JH_sJ^{-1}=-H_s$,
the real-linear operator $\mathcal A_s=H_sJ$
obeys $\mathcal A_s^2=H_s^2$.
For a computational string with
special-site bit $x_s$ and ordinary
Hamming weight $w$, the eigenvalue
of $H_s$ is
\[
 h(x)=\begin{cases}
 -2+w/2,&x_s=0,\\
 -1+w/2,&x_s=1.
 \end{cases}
\]
Hence $z=\mathcal A_sz$ implies $H_s^2z=z$,
restricting support to $h(x)^2=1$:
exactly $15+1+1+15=32$ strings, those of
ordinary weights $2,6$ at $x_s=0$ and $0,4$ at $x_s=1$.

On this shell the equation identifies the
weight-four amplitudes with the
\emph{negative complex conjugates}
of the complementary weight-two amplitudes,
and makes the two endpoint amplitudes complex conjugates.
A diagonal $SU(2)$ phase at the special site
makes both endpoint amplitudes real and
nonnegative. This is exactly the complex
pair/complement ansatz, and the converse
follows by the same spin-flip sign calculation.
Theorem~\ref{thm:complexsos} and
Proposition~\ref{prop:confphase} give the purity
assertions. The icosahedral relation follows
coefficientwise from~\eqref{eq:ico}.
\end{proof}
